\documentclass[11pt]{article}

\usepackage[margin=1in]{geometry}
\usepackage{amsmath,amssymb}
\usepackage{booktabs}
\usepackage{array}
\usepackage[hidelinks]{hyperref}

\newcommand{\Rhat}{\hat{R}}
\newcommand{\Aobs}{A}
\newcommand{\Amod}{\widehat{A}}
\newcommand{\Jatten}{J_{\mathrm{atten}}}
\newcommand{\Jreg}{J_{\mathrm{reg}}}
\newcommand{\argmin}{\operatorname*{arg\,min}}

\title{Previous Attenuation-Based Inverse Solvers for Rainfall Reconstruction}
\author{Working note}
\date{\today}

\begin{document}
\maketitle

\section*{Purpose}
This note summarizes earlier inverse solvers that reconstruct a rainfall field from microwave-link attenuation measurements.
The notation follows the notation used in \texttt{Journal.tex} so that the methods can be compared with our formulation.
The emphasis is on whether the method fits the link attenuation directly, rather than first reducing each link to a virtual point rain gauge and then interpolating.

\section*{Common notation}
Let $P$ denote the set of pixels in a reconstruction patch and let $\mathcal{L}$ denote the set of microwave links.
For a link $\ell\in\mathcal{L}$, let $A_\ell$ denote the observed rain-induced attenuation, let $|\ell|$ denote the link length, and let $l_{\ell p}$ denote the length of the part of link $\ell$ that crosses pixel $p$.
For an estimated rain-rate field $\Rhat:P\to\mathbb{R}_{\geq 0}$, the attenuation implied by the forward model is
\[
  A_\ell(\Rhat)
  \approx
  \sum_{p\in P} l_{\ell p}\,\kappa_\ell\,\Rhat(p)^{\alpha_\ell} .
\]
In our terminology, a direct attenuation-fitting method contains a term of the form
\[
  \Jatten(\Rhat)
  =
  \sum_{\ell\in\mathcal{L}} w_\ell\,
  \bigl(A_\ell(\Rhat)-A_\ell\bigr)^2,
\]
possibly after linearization, a change of variables, or replacement by an equivalent system of tomographic equations.
The main difference between the previous methods is how they regularize this ill-posed inverse problem.

\section{Giuli et al. (1991): microwave attenuation tomography}
\textbf{Reference.}
D. Giuli, A. Toccafondi, G. Biffi Gentili, and A. Freni,
\emph{Tomographic reconstruction of rainfall fields through microwave attenuation measurements},
Journal of Applied Meteorology, 1991.
DOI: \href{https://doi.org/10.1175/1520-0450(1991)030%3C1323:TRORFT%3E2.0.CO;2}{10.1175/1520-0450(1991)030<1323:TRORFT>2.0.CO;2}.

\paragraph{Decision variable in our notation.}
We can describe the Giuli et al. method using the same object as in our manuscript: an estimated rainfall field $\Rhat:P\to\mathbb{R}_{\geq 0}$.
The paper implements the reconstruction with smooth basis functions, but for comparison it is enough to regard $\Rhat(p)$ as the unknown value assigned to pixel or cell $p$.
In the frequency regime considered by the tomographic model, specific attenuation is treated as a linear or nearly linear function of rainfall.
Thus the unknown basis coefficients in the paper correspond, in our notation, to a smoothed representation of $\Rhat$ or of the induced specific attenuation $\kappa\Rhat$.

\paragraph{Forward model in our notation.}
For every link $\ell$, the model predicts the link attenuation from the reconstructed field by a path integral,
\[
  A_\ell(\Rhat)
  \approx
  \sum_{p\in P} l_{\ell p}\,\kappa\,\Rhat(p),
\]
where the linear coefficient $\kappa$ is fixed by the selected microwave frequency and polarization.
Equivalently, if a non-linear power law is retained, the same notation becomes
\[
  A_\ell(\Rhat)
  \approx
  \sum_{p\in P} l_{\ell p}\,\kappa\,\Rhat(p)^\alpha .
\]
For the original Giuli et al. tomography, the important point is that each measurement constrains an integral along the link path, not a point rain value.

\paragraph{Optimization problem in our notation.}
Using the notation of \texttt{Journal.tex}, the Giuli et al. reconstruction can be viewed as a regularized attenuation-fitting problem,
\[
  \min_{\Rhat}
  \underbrace{\sum_{\ell\in\mathcal{L}}
  \bigl(A_\ell(\Rhat)-A_\ell\bigr)^2}_{\Jatten(\Rhat)}
  +
  \lambda\,\Jreg(\Rhat).
\]
The regularization term $\Jreg$ is not identical to our $J_{1d}+J_{2d}+J_{\mathrm{total}}$.
It is induced by the tomographic reconstruction scheme and by the smooth basis representation.
In the linear basis implementation, this objective is equivalent to fitting a linear system for the link attenuations with a smoothness penalty.
Written in our language, however, the method is still minimizing disagreement between observed link attenuations $A_\ell$ and forward-modeled attenuations $A_\ell(\Rhat)$.

\paragraph{Comparison with our formulation.}
This is a direct precursor of attenuation-consistent inversion.
The main differences are that Giuli et al. considered a dedicated tomographic network and a linear or nearly linear attenuation representation.
Our formulation estimates pixel rain rates directly, uses the link-specific ITU-R forward model, imposes nonnegativity explicitly, and uses the explicit terms $\Jatten$, $J_{1d}$, $J_{2d}$, and $J_{\mathrm{total}}$.
The conceptual overlap is the attenuation-mismatch term.

\section{Zinevich, Alpert, and Messer (2008): nonlinear tomography on a variable-density grid}
\textbf{Reference.}
A. Zinevich, P. Alpert, and H. Messer,
\emph{Estimation of rainfall fields using commercial microwave communication networks of variable density},
Advances in Water Resources, 2008.
DOI: \href{https://doi.org/10.1016/j.advwatres.2008.03.003}{10.1016/j.advwatres.2008.03.003}.

\paragraph{Decision variable.}
The unknown is a rainfall vector $r=(r_1,
\ldots,r_n)$ on a variable-density reconstruction grid.
The grid is adapted to the commercial link topology, with smaller cells where links are denser.

\paragraph{Forward model.}
For link $i$, with power-law coefficients $a_i,b_i$, the observed attenuation is modeled as
\[
  A_i
  =
  a_i\int_{\ell_i} r(x)^{b_i}\,dx
  \approx
  a_i\sum_{j=1}^n l_{ij}r_j^{b_i} .
\]
Equivalently, if $R_i$ is the path-average rain rate inferred from the observed attenuation, the tomographic residuals are written as
\[
  f_i(r)=\sum_{j=1}^n l_{ij}r_j^{b_i}-L_iR_i^{b_i},
  \qquad i=1,\ldots,m .
\]
These residuals are attenuation-consistency equations, expressed after dividing by the link-specific power-law coefficient.

\paragraph{Optimization problem.}
The method is not written as a single scalar objective like our $J(\Rhat)$.
Instead, it solves the nonlinear system $f(r)=0$ by Newton-Raphson linearization.
At iteration $t$, the update $\delta r$ solves
\[
  J_t\,\delta r = -f(r_t),
\]
where $J_t$ is the Jacobian of the attenuation residuals.
Because the inverse problem is underdetermined, the linearized problem is solved by SIRT with smoothing and positivity constraints.
The smoothing is imposed through a spatial correlation operator derived from a rainfall correlation model.

\paragraph{Comparison with our formulation.}
This paper is very close in spirit to our attenuation term because it fits the path-integrated attenuation equations directly.
The important differences are methodological.
It solves nonlinear equations by repeated linearization and smoothing-constrained tomography, while our solver minimizes an explicit weighted objective containing $\Jatten$, first differences, second differences, and a total-rain shrinkage term.
Zinevich et al. also build the reconstruction grid around link density, whereas our benchmark uses a fixed pixel grid inherited from the rainfall patches.

\section{Bianchi et al. (2013): variational assimilation of radar, gauges, and microwave links}
\textbf{Reference.}
B. Bianchi, P. J. van Leeuwen, R. J. Hogan, and A. Berne,
\emph{A variational approach to retrieve rain rate by combining information from rain gauges, radars, and microwave links},
Journal of Hydrometeorology, 2013.
DOI: \href{https://doi.org/10.1175/JHM-D-12-094.1}{10.1175/JHM-D-12-094.1}.

\paragraph{Decision variable.}
The unknown is a state vector $x$ that includes rain-rate variables over pixels and additional calibration parameters.
A radar rain field provides a background estimate $x_b$.
Rain gauges and microwave links are used as observations that update this background.

\paragraph{Forward model.}
The observation operator $H(x)$ maps the rain-rate state to the observations.
For microwave links, the link path is divided into segments crossing pixels, and the path-averaged or total link quantity is computed from the pixel values and segment lengths.
The microwave-link part of $H(x)$ is therefore a path-integrated observation operator, not a point-gauge interpolation operator.

\paragraph{Optimization problem.}
The variational retrieval minimizes a data-assimilation cost function of the form
\[
  \min_x
  \frac{1}{2}(x-x_b)^\top B^{-1}(x-x_b)
  +
  \frac{1}{2}\bigl(y-H(x)\bigr)^\top R^{-1}\bigl(y-H(x)\bigr),
\]
where $B$ is the background-error covariance matrix and $R$ is the observation-error covariance matrix.
The observation vector $y$ contains rain-gauge and microwave-link observations.
The microwave-link term is an attenuation-related mismatch term weighted by the estimated link-observation uncertainty.
The cost function is minimized using a Gauss-Newton method.

\paragraph{Comparison with our formulation.}
This is a variational inverse problem with an explicit observation-mismatch term.
Bianchi et al. solve a multisensor assimilation problem with radar as a prior field. In contrast, our benchmark isolates the CML inverse problem and uses no radar prior in the solver.
Their regularization enters mainly through the background term and covariance matrices.
Our regularization is expressed directly as spatial smoothness and total-rain terms on the reconstructed field.

\section{Roy, Gishkori, and Leus (2016): dynamic sparse rainfall monitoring}
\textbf{Reference.}
V. Roy, S. Gishkori, and G. Leus,
\emph{Dynamic rainfall monitoring using microwave links},
EURASIP Journal on Advances in Signal Processing, 2016.
DOI: \href{https://doi.org/10.1186/s13634-016-0367-6}{10.1186/s13634-016-0367-6}.

\paragraph{Decision variable.}
The unknown is a time-dependent rainfall vector $u_t\in\mathbb{R}^N$ over a fixed grid.
The method also uses a sparse representation $u_t=\Psi_t z_t$.

\paragraph{Forward model.}
The nonlinear link-attenuation model is
\[
  y_{i,t}
  =
  a\sum_{j=1}^N l_{ij}u_{j,t}^{b} + e_{i,t},
  \qquad i=1,\ldots,M .
\]
In vector form,
\[
  y_t=\Phi(u_t)+e_t .
\]
The rainfall dynamics are modeled by
\[
  u_t=H_tu_{t-1}+q_t .
\]

\paragraph{Optimization problem.}
At each time step, the nonlinear observation map is linearized around the predicted field $\hat{u}_{t|t-1}$.
The correction step solves the constrained sparse optimization problem
\[
  \min_{z_t:\,\Psi_t z_t\geq 0}
  \left\|\hat{u}_{t|t-1}-\Psi_t z_t\right\|^2_{M_{t|t-1}^{-1}}
  +
  \left\|\tilde{y}_t-J_t\Psi_t z_t\right\|^2_{R^{-1}}
  +
  \lambda_t\|z_t\|_1 .
\]
Here $J_t$ is the Jacobian of the attenuation observation map and $\tilde{y}_t$ is the linearized observation vector.
The first term penalizes disagreement with the dynamic prediction, the second term penalizes attenuation mismatch after linearization, and the third term enforces sparsity.
The constraint imposes nonnegative rainfall.

\paragraph{Comparison with our formulation.}
This is one of the closest formulations to our constrained optimization language.
It explicitly combines attenuation mismatch, nonnegativity, and additional prior structure.
The main difference is that Roy et al. solve a dynamic state-estimation problem and use a sparsity penalty in a basis representation, while our current solver is static for each benchmark patch and uses first-difference, second-difference, and total-rain penalties.

\section{D'Amico, Manzoni, and Solazzi (2016): operational CML tomography}
\textbf{Reference.}
M. D'Amico, A. Manzoni, and G. L. Solazzi,
\emph{Use of operational microwave link measurements for the tomographic reconstruction of 2-D maps of accumulated rainfall},
IEEE Geoscience and Remote Sensing Letters, 2016.
DOI: \href{https://doi.org/10.1109/LGRS.2016.2614326}{10.1109/LGRS.2016.2614326}.

\paragraph{Decision variable.}
The unknown is again a basis expansion of the specific attenuation field,
\[
  \hat{\gamma}(x,y)=\sum_{n=1}^N \gamma_n b_n(x,y).
\]
The estimated specific attenuation is converted to rainfall accumulation.

\paragraph{Forward model.}
The link measurements form the same linear tomographic system used in the Giuli-type approach,
\[
  p=A\gamma+e,
\]
where $p$ contains measured link attenuations and $A_{mn}$ is obtained by integrating basis function $b_n$ along link $m$.

\paragraph{Optimization problem.}
The paper applies the recursive tomographic reconstruction algorithm
\[
  \gamma^{(l+1)}
  =
  \gamma^{(l)}
  +
  \frac{1}{k}
  \left[-(I+b_0B)\gamma^{(l)} + a_0A^\top\bigl(p-A\gamma^{(l)}\bigr)\right].
\]
This is equivalent to an iterative descent method for a regularized attenuation-mismatch problem of the form
\[
  \min_\gamma
  \frac{a_0}{2}\|p-A\gamma\|_2^2
  +
  \frac{1}{2}\gamma^\top(I+b_0B)\gamma .
\]
The parameters $a_0$ and $b_0$ are calibrated using rainy days for which gauge information is available.

\paragraph{Comparison with our formulation.}
D'Amico et al. demonstrate that attenuation-based tomography can be applied to operational links, not only to simulated dedicated networks.
The method fits link attenuation directly, but it reconstructs specific attenuation through a linear tomographic basis expansion.
Our method reconstructs the rain-rate field directly on the benchmark pixel grid and treats the ITU-R power law inside the attenuation term.

\section*{Summary table}
\begin{center}
\small
\renewcommand{\arraystretch}{1.25}
\begin{tabular}{>{\raggedright\arraybackslash}p{0.20\textwidth}>{\raggedright\arraybackslash}p{0.21\textwidth}>{\raggedright\arraybackslash}p{0.22\textwidth}>{\raggedright\arraybackslash}p{0.25\textwidth}}
\toprule
Paper & Decision variable & Attenuation fit & Regularization or prior \\
\midrule
Giuli et al. (1991) & Rain field $\Rhat$ represented by smooth basis functions & Linearized $A_\ell(\Rhat)-A_\ell$ residual & Recursive regularized tomography \\
Zinevich et al. (2008) & Rain-rate vector $r$ on a variable-density grid & Nonlinear residuals $f_i(r)=0$ from link attenuation equations & SIRT, smoothing operator, positivity \\
Bianchi et al. (2013) & Rain-rate state plus calibration parameters & Microwave-link observation term in a variational cost & Radar background covariance and observation-error covariance \\
Roy et al. (2016) & Dynamic rain-rate vector $u_t=\Psi_tz_t$ & Linearized attenuation residual in each correction step & Dynamic prediction, sparsity, nonnegativity \\
D'Amico et al. (2016) & Field represented by basis coefficients & Linearized $A_\ell(\Rhat)-A_\ell$ residual & Recursive tomography with calibrated regularization \\
\bottomrule
\end{tabular}
\end{center}

\section*{Takeaway for our manuscript}
It is not accurate to claim that attenuation-mismatch inversion is new.
The safer claim is that our work revisits attenuation-consistent reconstruction in a controlled benchmark and compares it directly with IDW-type virtual-gauge baselines.
A possible manuscript sentence is:

\begin{quote}
Previous tomographic and variational methods have reconstructed rainfall fields by fitting path-integrated microwave attenuation measurements directly.
Our contribution is to place this attenuation-consistent inverse formulation in a controlled noise-free benchmark, to compare it systematically with IDW-type virtual-gauge baselines, and to study the effect of link geometry through distance-to-link stratification.
\end{quote}

\end{document}
